\documentclass[10pt, a4paper]{article}
\usepackage{graphicx}

\title{T22 experiments report, \textit{Preliminary}}
\author{Leios Red team}
\date{28.09.2026, ver. 1.0}

\begin{document}
\maketitle

\subsection*{Experiment set-up}

In order to test possibilities in T20 attack, a local Piranha cluster was 
established, and tested under various (adversarial) test scenarios.

Some nodes are ``honest'' and do not participate in adversarial activity,
others are ``adversarial''. Typical scenario is a combination of normal work 
(all nodes behave honestly) and an attack (adversarial part performs an attack).
Typical layout is 10+10+10 mins run (10 min. normal, 10 min. 
adversarial, 10 min. normal).

The default parameters for the network are the following:

\begin{itemize}
\item Piranha (net-rs) network, 30 Piranha nodes.
\item Basic topology: random, 30 nodes, network diameter 4.
\item Basic txs tempo: transactions generated uniformly by all nodes,
1 tx per node per second.
\end{itemize}

The attacks are presented roughly in the power growth proportion. 
Ultimate target of the attack: stop EB certification, and it's partially 
achieved.

%The variation in protocol implementation gives us some space to research 
%general protocol performance. The second implementation was introduced after 
%the first one demonstrated unreliable behaviour.

\section*{Raw results}


\end{document}
